% GENERATED FILE. Edit canonical lessons, then run npm run generate:books.
% canonical-source-hash: fnv1a64:4ecc864c265308cc
% canonical-lessons: ML-C379-alkkuttam, ML-C379-itikkuka, ML-C379-kuttimuttuka, ML-C379-mariyuka, ML-C379-tattiviluka

\chapter{A Crash at the Junction}
\label{ch:ml-a-crash-at-the-junction}
\hlchaptermodality{379}

\begin{chapteropening}
\textbf{By the end of this chapter:} \emph{I can say a throng, to crash into, to collide, to overturn and to trip and fall.}

Complete the last lesson of chapter 379: I can say a throng, to crash into, to collide, to overturn and to trip and fall.
\end{chapteropening}

\section[\texorpdfstring{\ml{ആൾക്കൂട്ടം} (āḷkkūṭṭaṁ)}{ആൾക്കൂട്ടം (āḷkkūṭṭaṁ)}]{\ml{ആൾക്കൂട്ടം} (āḷkkūṭṭaṁ) — a throng; a crowd of people}

\label{lesson:ML-C379-alkkuttam}



\begin{quote}
Before the new one: say the Malayalam for to flee, then the Malayalam for to seek.
\end{quote}

\subsection*{You'll want to know: \ml{ആൾക്കൂട്ടം}}
\textbf{\ml{ആൾക്കൂട്ടം}} — \emph{āḷkkūṭṭaṁ} — \textquotedblleft{}a throng; a crowd of people\textquotedblright{}.

\begin{quote}
\ml{അപകടം} \ml{നടന്ന} \ml{സ്ഥലത്ത്} \ml{വലിയ} \ml{ആൾക്കൂട്ടം} \ml{ഉണ്ടായിരുന്നു}.
\end{quote}

\begin{quote}
\emph{apakaṭaṁ naṭanna sthalattŭ valiya āḷkkūṭṭaṁ uṇṭāyirunnu.}
\end{quote}

\begin{quote}
There was a big crowd where the accident happened.
\end{quote}

\ml{ആൾ}, a person, + \ml{കൂട്ടം}, a group.

\begin{cousinweb}
\textbf{\ml{ആൾക്കൂട്ടം}} (\emph{āḷkkūṭṭaṁ}) is a native Dravidian compound of \ml{ആൾ} (\emph{āḷ}), \textquotedblleft{}a person\textquotedblright{}, and \ml{കൂട്ടം}, \textquotedblleft{}a gathering, a group\textquotedblright{}, from the root of \ml{കൂടുക}, \textquotedblleft{}to come together\textquotedblright{} (DEDR 1882).
\end{cousinweb}

\begin{grammarlens}[title={What you've built}]
One more word for telling what happened.
\end{grammarlens}

\subsection*{Your turn}
\begin{itemize}
\raggedright
  \item \emph{Say it:} \emph{āḷkkūṭṭaṁ}
  \item \emph{Say it:} \emph{āḷkkūṭṭaṁ}, once more
  \item \emph{Say it:} \emph{apakaṭaṁ naṭanna sthalattŭ valiya āḷkkūṭṭaṁ uṇṭāyirunnu.} — There was a big crowd where the accident happened.
  \item \emph{Recall:} say the Malayalam for to flee, then the Malayalam for to seek, then say \emph{āḷkkūṭṭaṁ} again
\end{itemize}

\begin{tcolorbox}[breakable,colback=teal!4,colframe=teal!35!black,title={Before you move on}]
What does \ml{ആൾക്കൂട്ടം} mean? (\textquotedblleft{}A throng; a crowd of people\textquotedblright{}.) Say it once more.
\end{tcolorbox}
\section[\texorpdfstring{\ml{ഇടിക്കുക} (iṭikkuka)}{ഇടിക്കുക (iṭikkuka)}]{\ml{ഇടിക്കുക} (iṭikkuka) — to crash into; to hit}

\label{lesson:ML-C379-itikkuka}



\begin{quote}
Before the new one: say the Malayalam for to seek, then the Malayalam for a throng.
\end{quote}

\subsection*{You'll want to know: \ml{ഇടിക്കുക}}
\textbf{\ml{ഇടിക്കുക}} — \emph{iṭikkuka} — \textquotedblleft{}to crash into; to hit\textquotedblright{}.

\begin{quote}
\ml{ഒരു} \ml{കാർ} \ml{ബസിൽ} \ml{ഇടിച്ചു}.
\end{quote}

\begin{quote}
\emph{oru kār basil iṭiccu.}
\end{quote}

\begin{quote}
A car crashed into a bus.
\end{quote}

Past: \ml{ഇടിച്ചു}, \emph{iṭiccu}.

What is hit takes -\ml{ഇൽ}: \ml{മരത്തിൽ} \ml{ഇടിച്ചു}, hit a tree.

\begin{cousinweb}
\textbf{\ml{ഇടിക്കുക}} (\emph{iṭikkuka}) is native Dravidian: DEDR 443 lists Malayalam \emph{iṭikka}, \textquotedblleft{}to beat, to bruise\textquotedblright{}, and \emph{iṭi}, \textquotedblleft{}a stroke, a blow, a shock\textquotedblright{}.
\end{cousinweb}

\begin{grammarlens}[title={What you've built}]
One more word for telling what happened.
\end{grammarlens}

\subsection*{Your turn}
\begin{itemize}
\raggedright
  \item \emph{Say it:} \emph{iṭikkuka}
  \item \emph{Say it:} \emph{iṭikkuka}, once more
  \item \emph{Say it:} \emph{oru kār basil iṭiccu.} — A car crashed into a bus.
  \item \emph{Recall:} say the Malayalam for to seek, then the Malayalam for a throng, then say \emph{iṭikkuka} again
\end{itemize}

\begin{tcolorbox}[breakable,colback=teal!4,colframe=teal!35!black,title={Before you move on}]
What does \ml{ഇടിക്കുക} mean? (\textquotedblleft{}To crash into; to hit\textquotedblright{}.) Say it once more.
\end{tcolorbox}
\section[\texorpdfstring{\ml{കൂട്ടിമുട്ടുക} (kūṭṭimuṭṭuka)}{കൂട്ടിമുട്ടുക (kūṭṭimuṭṭuka)}]{\ml{കൂട്ടിമുട്ടുക} (kūṭṭimuṭṭuka) — to collide}

\label{lesson:ML-C379-kuttimuttuka}



\begin{quote}
Before the new one: say the Malayalam for a throng, then the Malayalam for to crash into.
\end{quote}

\subsection*{You'll want to know: \ml{കൂട്ടിമുട്ടുക}}
\textbf{\ml{കൂട്ടിമുട്ടുക}} — \emph{kūṭṭimuṭṭuka} — \textquotedblleft{}to collide\textquotedblright{}.

\begin{quote}
\ml{കവലയിൽ} \ml{രണ്ട്} \ml{ബൈക്കുകൾ} \ml{കൂട്ടിമുട്ടി}.
\end{quote}

\begin{quote}
\emph{kavalayil raṇṭŭ baikkukaḷ kūṭṭimuṭṭi.}
\end{quote}

\begin{quote}
Two motorbikes collided at the junction.
\end{quote}

Past: \ml{കൂട്ടിമുട്ടി}, \emph{kūṭṭimuṭṭi}.

Both sides hit each other: \ml{കൂട്ടി}, together, + \ml{മുട്ടുക}, to knock.

\begin{cousinweb}
\textbf{\ml{കൂട്ടിമുട്ടുക}} (\emph{kūṭṭimuṭṭuka}) is a native Dravidian compound of \ml{കൂട്ടി}, \textquotedblleft{}together\textquotedblright{}, from \ml{കൂടുക} (DEDR 1882), and \ml{മുട്ടുക} (\emph{muṭṭuka}), \textquotedblleft{}to dash against, to knock\textquotedblright{} (DEDR 4932).
\end{cousinweb}

\begin{grammarlens}[title={What you've built}]
One more word for telling what happened.
\end{grammarlens}

\subsection*{Your turn}
\begin{itemize}
\raggedright
  \item \emph{Say it:} \emph{kūṭṭimuṭṭuka}
  \item \emph{Say it:} \emph{kūṭṭimuṭṭuka}, once more
  \item \emph{Say it:} \emph{kavalayil raṇṭŭ baikkukaḷ kūṭṭimuṭṭi.} — Two motorbikes collided at the junction.
  \item \emph{Recall:} say the Malayalam for a throng, then the Malayalam for to crash into, then say \emph{kūṭṭimuṭṭuka} again
\end{itemize}

\begin{tcolorbox}[breakable,colback=teal!4,colframe=teal!35!black,title={Before you move on}]
What does \ml{കൂട്ടിമുട്ടുക} mean? (\textquotedblleft{}To collide\textquotedblright{}.) Say it once more.
\end{tcolorbox}
\section[\texorpdfstring{\ml{മറിയുക} (maṟiyuka)}{മറിയുക (maṟiyuka)}]{\ml{മറിയുക} (maṟiyuka) — to overturn; to tip over}

\label{lesson:ML-C379-mariyuka}



\begin{quote}
Before the new one: say the Malayalam for to crash into, then the Malayalam for to collide.
\end{quote}

\subsection*{You'll want to know: \ml{മറിയുക}}
\textbf{\ml{മറിയുക}} — \emph{maṟiyuka} — \textquotedblleft{}to overturn; to tip over\textquotedblright{}.

\begin{quote}
\ml{വളവിൽ} \ml{ബസ്} \ml{മറിഞ്ഞു}.
\end{quote}

\begin{quote}
\emph{vaḷavil basŭ maṟiññu.}
\end{quote}

\begin{quote}
The bus overturned on the bend.
\end{quote}

Past: \ml{മറിഞ്ഞു}, \emph{maṟiññu}.

Of vehicles, boats and pots.

\begin{cousinweb}
\textbf{\ml{മറിയുക}} (\emph{maṟiyuka}) is native Dravidian: DEDR 4761 lists Malayalam \emph{maṟiyuka}, \textquotedblleft{}to turn back, to turn over, to be upset\textquotedblright{}.
\end{cousinweb}

\begin{grammarlens}[title={What you've built}]
One more word for telling what happened.
\end{grammarlens}

\subsection*{Your turn}
\begin{itemize}
\raggedright
  \item \emph{Say it:} \emph{maṟiyuka}
  \item \emph{Say it:} \emph{maṟiyuka}, once more
  \item \emph{Say it:} \emph{vaḷavil basŭ maṟiññu.} — The bus overturned on the bend.
  \item \emph{Recall:} say the Malayalam for to crash into, then the Malayalam for to collide, then say \emph{maṟiyuka} again
\end{itemize}

\begin{tcolorbox}[breakable,colback=teal!4,colframe=teal!35!black,title={Before you move on}]
What does \ml{മറിയുക} mean? (\textquotedblleft{}To overturn; to tip over\textquotedblright{}.) Say it once more.
\end{tcolorbox}
\section[\texorpdfstring{\ml{തട്ടിവീഴുക} (taṭṭivīḻuka)}{തട്ടിവീഴുക (taṭṭivīḻuka)}]{\ml{തട്ടിവീഴുക} (taṭṭivīḻuka) — to trip and fall}

\label{lesson:ML-C379-tattiviluka}



\begin{quote}
Before the new one: say the Malayalam for to collide, then the Malayalam for to overturn.
\end{quote}

\subsection*{You'll want to know: \ml{തട്ടിവീഴുക}}
\textbf{\ml{തട്ടിവീഴുക}} — \emph{taṭṭivīḻuka} — \textquotedblleft{}to trip and fall\textquotedblright{}.

\begin{quote}
\ml{അമ്മ} \ml{പടിയിൽ} \ml{തട്ടിവീണു}.
\end{quote}

\begin{quote}
\emph{amma paṭiyil taṭṭivīṇu.}
\end{quote}

\begin{quote}
Mother tripped on the step and fell.
\end{quote}

Past: \ml{തട്ടിവീണു}, \emph{taṭṭivīṇu}.

\ml{തട്ടുക}, to knock against, + \ml{വീഴുക}, to fall.

\begin{cousinweb}
\textbf{\ml{തട്ടിവീഴുക}} (\emph{taṭṭivīḻuka}) is a native Dravidian compound of \ml{തട്ടുക} (\emph{taṭṭuka}), \textquotedblleft{}to strike against, to knock\textquotedblright{} (DEDR 3039), and \ml{വീഴുക} (\emph{vīḻuka}), \textquotedblleft{}to fall\textquotedblright{}.
\end{cousinweb}

\begin{grammarlens}[title={What you've built}]
One more word for telling what happened.
\end{grammarlens}

\subsection*{Your turn}
\begin{itemize}
\raggedright
  \item \emph{Say it:} \emph{taṭṭivīḻuka}
  \item \emph{Say it:} \emph{taṭṭivīḻuka}, once more
  \item \emph{Say it:} \emph{amma paṭiyil taṭṭivīṇu.} — Mother tripped on the step and fell.
  \item \emph{Recall:} say the Malayalam for to collide, then the Malayalam for to overturn, then say \emph{taṭṭivīḻuka} again
\end{itemize}

\begin{tcolorbox}[breakable,colback=teal!4,colframe=teal!35!black,title={Before you move on}]
What does \ml{തട്ടിവീഴുക} mean? (\textquotedblleft{}To trip and fall\textquotedblright{}.) Say it once more.
\end{tcolorbox}
